%%
%% part1-getting-started/ch01-introduction.tex
%%
%% Chapter 1: Introduction to MatinBook — What it is, what it
%% can do, and how it compares to other LaTeX classes.
%%

\chapter{معرفی و قابلیت‌ها}
\label{chap:introduction}

\section{چرا این موضوع مهم است؟}
\label{sec:introduction-why}

اگر تا به حال تلاش کرده‌اید یک کتاب فنی فارسی با
ابزارهای حرفه‌ای حروف‌چینی کنید، احتمالاً با
چالش‌های متعددی روبه‌رو شده‌اید: فونت‌های فارسی
که در اندازه‌های کوچک خوانا نیستند، فرمول‌های
ریاضی که با متن فارسی قاطی می‌شوند، کدهای
برنامه‌نویسی که جهت متن را به‌هم می‌ریزند، و
ده‌ها مشکل کوچک و بزرگ دیگر.

\lr{MatinBook} برای حل همین مشکلات طراحی شده است.
این کلاس \lr{LaTeX}، تمام ابزارهای لازم برای
تولید یک کتاب فنی فارسی حرفه‌ای را در یک بسته‌ی
منسجم فراهم می‌کند. هدف این فصل، آشنایی اولیه با
\lr{MatinBook}، قابلیت‌های کلیدی آن، و تفاوت‌هایش
با سایر ابزارهای موجود است.

\mnote{\lr{MatinBook} یک کلاس \lr{LaTeX} است، نه یک بسته‌ی ساده. این یعنی تمام قابلیت‌های کتاب‌سازی در آن یکپارچه شده است.}

\mnote{نام \lr{MatinBook} از نام کوچک توسعه‌دهنده‌ی آن، «متین»، گرفته شده است.}

\section{\lr{MatinBook} چیست؟}
\label{sec:introduction-what}

\lr{MatinBook} یک کلاس \lr{LaTeX} است که به‌طور
اختصاصی برای نوشتن کتاب‌های فنی و علمی فارسی طراحی
شده است. این کلاس توسط متین محمودی توسعه یافته و
تحت مجوز \lr{MIT} منتشر می‌شود.

\begin{definition}[\lr{MatinBook}]
    \lr{MatinBook} یک کلاس \lr{LaTeX} است که با
    موتور \lr{XeLaTeX} کامپایل می‌شود و تمام
    ابزارهای لازم برای تولید کتاب‌های فنی فارسی
    — از حروف‌چینی راست‌به‌چپ تا نمودارهای پیچیده —
    را در یک ساختار ماژولار فراهم می‌کند.
    \label{def:matinbook}
\end{definition}

\lr{MatinBook} بر پایه‌ی کلاس استاندارد \lr{book}
ساخته شده است، ولی قابلیت‌های زیادی به آن اضافه
کرده است:

\begin{itemize}
    \item پشتیبانی کامل از حروف‌چینی راست‌به‌چپ
    فارسی با \lr{XePersian}.
    \item محیط‌های علمی رنگ‌آمیزی‌شده (قضیه، تعریف،
    مثال، اثبات).
    \item پشتیبانی از کدهای برنامه‌نویسی با
    رنگ‌آمیزی نحوی و کامنت فارسی.
    \item نمودارها و شکل‌های حرفه‌ای با \lr{TikZ} و
    \lr{PGFPlots}.
    \item جداول مدرن با \lr{tabularray}.
    \item مدیریت منابع با \lr{biblatex} و \lr{biber}.
    \item نمایه‌ی حرفه‌ای با \lr{xindy} و مرتب‌سازی
    صحیح فارسی.
    \item جلد تمام‌رنگی با طرح حرفه‌ای.
    \item یادداشت‌های حاشیه‌ای برای توضیحات مکمل.
\end{itemize}

\mnote{کلمه‌ی کلیدی اینجا «یکپارچه» است. در ابزارهای دیگر، باید هر قابلیت را دستی پیکربندی کنید.}

\section{قابلیت‌های کلیدی}
\label{sec:introduction-features}

در این بخش، نگاهی گذرا به مهم‌ترین قابلیت‌های
\lr{MatinBook} می‌اندازیم. هر یک از این قابلیت‌ها
در فصل‌های بعدی به‌طور مفصل بررسی خواهد شد.

\subsection{محیط‌های علمی رنگ‌آمیزی‌شده}
\label{sec:introduction-environments}

\lr{MatinBook} ده نوع محیط علمی را با دو خانواده‌ی
بصری متمایز ارائه می‌دهد. محیط‌های نظری (قضیه، لم،
نتیجه، گزاره) از یک نوار رنگی توپر استفاده
می‌کنند، در حالی که محیط‌های توضیحی (مثال، اثبات،
راه‌حل) از یک خط عمودی سمت راست بهره می‌برند.

\mnote{تفکیک دو خانواده، بر اساس نوع محتوا است، نه سلیقه. محیط‌های نظری «رسمی» هستند، محیط‌های توضیحی «غیررسمی».}

\begin{theorem}[قضیه‌ی فیثاغورث]
    در هر مثلث قائم‌الزاویه، مربع وتر برابر است با
    مجموع مربعات دو ضلع دیگر:
    \[
        a^{2} + b^{2} = c^{2}
    \]
    \label{thm:pythagoras-intro}
\end{theorem}

\begin{proof}
    اثبات این قضیه را در فصل \cref{chap:environments}
    به‌طور کامل بررسی خواهیم کرد.
\end{proof}

\begin{example}
    برای یک مثلث با ضلع‌های $a = 3$ و $b = 4$،
    وتر برابر است با:
    \[
        c = \sqrt{3^{2} + 4^{2}} = \sqrt{25} = 5
    \]
    \label{ex:pythagoras-intro}
\end{example}

\begin{remark}
    محیط‌های علمی در \lr{MatinBook} از یک شمارنده‌ی
    مشترک استفاده می‌کنند. یعنی اگر یک تعریف با
    شماره‌ی ۱.۱ ثبت شود، قضیه‌ی بعدی شماره‌ی ۱.۲
    خواهد داشت.
    \label{rem:shared-counter-intro}
\end{remark}

این محیط‌ها در فصل \cref{chap:environments} به‌طور
مفصل بررسی می‌شوند.

\subsection{پشتیبانی از ریاضیات}
\label{sec:introduction-math}

\lr{MatinBook} از تمام قابلیت‌های \lr{amsmath}،
\lr{mathtools} و \lr{unicode-math} پشتیبانی می‌کند
و مجموعه‌ای از عملگرها و جداکننده‌های سفارشی را
نیز اضافه کرده است:

\begin{itemize}
    \item عملگرهای حسابان برداری: \cmd{grad}،
    \cmd{curl}، \cmd{diver}.
    \item عملگرهای جبر خطی: \cmd{rank}، \cmd{trace}،
    \cmd{diag}.
    \item مجموعه‌های اعداد: $\N$، $\Z$، $\Q$، $\R$،
    $\C$.
    \item جداکننده‌های هوشمند: \cmd{abs}، \cmd{norm}،
    \cmd{ceil}، \cmd{floor}.
\end{itemize}

مثال:

\[
    \grad f = \nabla f, \quad
    \diver \vec{F} = \nabla \cdot \vec{F}, \quad
    \norm{\vec{v}} = \sqrt{\inner{\vec{v}}{\vec{v}}}
\]

\mnote{در متن فارسی، فرمول‌های ریاضی باید داخل \lr{\textbackslash lr\{...\}} قرار گیرند تا جهت‌نمایی درست حفظ شود.}

ریاضیات در \lr{MatinBook} در فصل \cref{chap:math}
به‌طور کامل بررسی می‌شود.

\subsection{کد و الگوریتم}
\label{sec:introduction-code}

\lr{MatinBook} از بسته‌ی \lr{minted} برای نمایش کد
با رنگ‌آمیزی نحوی استفاده می‌کند. این بسته از بیش
از ۳۰۰ زبان برنامه‌نویسی پشتیبانی می‌کند و امکان
نوشتن کامنت فارسی داخل کد را نیز فراهم می‌سازد.

برای کامنت فارسی در پایتون، از دستور \cmd{pcc}
استفاده می‌کنیم که به‌طور خودکار علامت \lr{\#} را
قبل از متن فارسی قرار می‌دهد:

\begin{latin}
\begin{minted}{python}
def factorial(n: int) -> int:
    """Calculate n! recursively."""
    if n <= 1:
        return 1
    return n * factorial(n - 1)

\pcc{محاسبه‌ی فاکتوریل ۵}
print(factorial(5))  # 120
\end{minted}
\end{latin}
\codecaption{تابع فاکتوریل بازگشتی در پایتون}
\label{code:factorial-intro}

\mnote{برای کامنت فارسی در پایتون، \lr{\textbackslash pcc} بهترین گزینه است. برای زبان‌های دیگر (مثل \lr{LaTeX} یا \lr{bash}) می‌توانید از \lr{\textbackslash pc} داخل \lr{|...|} استفاده کنید.}

\mnote{اگر \lr{\textbackslash pcc} را دستی باز کنید، به \lr{|\textbackslash pc\{\textbackslash\#~متن\}|} بسط می‌یابد. این یعنی \lr{\#} داخل \lr{\textbackslash pc} قرار می‌گیرد تا \lr{Pygments} آن را به‌عنوان شروع کامنت نبیند.}

الگوریتم‌ها نیز با بسته‌ی استاندارد \lr{algorithm}
نمایش داده می‌شوند. نکته‌ی مهم این است که
\lr{xepersian} عنوان فارسی «الگوریتم» را به‌طور
خودکار اضافه می‌کند، ولی متن داخل الگوریتم باید
لاتین باشد (به‌جز کامنت‌هایی که با \cmd{pc} نوشته
می‌شوند):

\begin{latin}
\begin{algorithm}
\caption{جستجوی دودویی}
\label{alg:binary-intro}
\begin{algorithmic}[1]
    \Require Sorted array $A$ of $n$ elements
    \Require Target value $x$
    \Ensure Index of $x$ in $A$, or $-1$
    \State $left \gets 0$, $right \gets n - 1$
    \While{$left \leq right$}
        \State $mid \gets (left + right) / 2$
        \If{$A[mid] = x$}
            \State \Return $mid$  \Comment{|\pc{یافتن مقدار}|}
        \ElsIf{$A[mid] < x$}
            \State $left \gets mid + 1$  \Comment{|\pc{جستجو در نیمه‌ی راست}|}
        \Else
            \State $right \gets mid - 1$  \Comment{|\pc{جستجو در نیمه‌ی چپ}|}
        \EndIf
    \EndWhile
    \State \Return $-1$  \Comment{|\pc{یافت نشد}|}
\end{algorithmic}
\end{algorithm}
\end{latin}

\mnote{در الگوریتم‌ها، متن داخل \lr{\textbackslash Comment\{...\}} می‌تواند فارسی باشد، ولی باید با \lr{\textbackslash pc} پیچیده شود.}

کد و الگوریتم در \cref{chap:code} به‌طور مفصل
بررسی می‌شوند.

\subsection{نمودارها و شکل‌ها}
\label{sec:introduction-graphics}

\lr{MatinBook} از \lr{TikZ}، \lr{PGFPlots} و
\lr{pgf-pie} برای ترسیم نمودارها و شکل‌های حرفه‌ای
استفاده می‌کند. از نمودارهای ساده‌ی توابع تا
نمودارهای سه‌بعدی، فلوچارت‌ها، درخت‌ها و گراف‌ها،
همه در \lr{MatinBook} پشتیبانی می‌شوند.

\begin{figure}[h]
\centering
\begin{latin}
\begin{tikzpicture}
  \begin{axis}[
    width=0.8\textwidth, height=5cm,
    xlabel={$x$},
    ylabel={$f(x)$},
    grid=major,
    legend pos=north west,
    domain=-2:2,
    samples=100
  ]
    \addplot[matinblue, thick] {x^2};
    \addlegendentry{$x^2$}
    \addplot[matinred, thick] {x^3};
    \addlegendentry{$x^3$}
  \end{axis}
\end{tikzpicture}
\end{latin}
\caption{نمودار دو تابع $x^{2}$ و $x^{3}$}
\label{fig:functions-intro}
\end{figure}

\mnote{برای متن فارسی داخل \lr{tikzpicture}، باید از \lr{\textbackslash rl\{...\}} استفاده کنید.}

نمودار \cref{fig:functions-intro} نمونه‌ای ساده از
قابلیت‌های گرافیکی \lr{MatinBook} است. این موضوع
در \cref{chap:graphics} به‌طور کامل بررسی می‌شود.

\subsection{جداول مدرن}
\label{sec:introduction-tables}

\lr{MatinBook} از بسته‌ی \lr{tabularray} برای
جداول استفاده می‌کند. این بسته، برخلاف \lr{tabularx}،
با \lr{xepersian} کاملاً سازگار است و از
حروف‌چینی راست‌به‌چپ پشتیبانی می‌کند.

\begin{matintable}{مقایسه‌ی سه الگوریتم مرتب‌سازی}{tab:sorting-intro}
\begin{tblr}{
    width = \textwidth,
    colspec = {l X[c] X[c]},
    hlines, vlines,
    colsep = 8pt,
    rowsep = 4pt,
    row{1} = {font=\bfseries},
}
الگوریتم & بهترین حالت & بدترین حالت \\
مرتب‌سازی سریع & \lr{$O(n \log n)$} & \lr{$O(n^2)$} \\
مرتب‌سازی ادغامی & \lr{$O(n \log n)$} & \lr{$O(n \log n)$} \\
مرتب‌سازی حبابی & \lr{$O(n)$} & \lr{$O(n^2)$} \\
\end{tblr}
\end{matintable}

\mnote{بسته‌ی \lr{tabularx} با \lr{xepersian} ناسازگار است و خطا می‌دهد. همیشه از \lr{tabularray} استفاده کنید.}

جدول \cref{tab:sorting-intro} نمونه‌ای از جدول‌های
\lr{MatinBook} است. فصل \cref{chap:tables} به‌طور
کامل به این موضوع می‌پردازد.

\subsection{جلد حرفه‌ای}
\label{sec:introduction-cover}

\lr{MatinBook} یک جلد تمام‌رنگی با طرح حرفه‌ای ارائه
می‌دهد. جلد جلو شامل پس‌زمینه‌ی سرمه‌ای، گوه‌ی
قرمز، بیضی طلایی و نمونه‌کد \lr{C++} است. جلد
عقب شامل دو بلوک ایستاتیک — «درباره این کتاب» و
«درباره این نویسنده» — با فاصله‌ی دقیق
\pnum{۰.۷} سانتی‌متر است.

\begin{figure}[h]
\centering
\begin{latin}
\begin{tikzpicture}[scale=0.6]
  % Simplified front cover mockup
  \fill[matinnavy] (0,0) rectangle (6,8);
  \fill[matincrimson] (0,8) -- (3.5,8) -- (2,6.5) -- cycle;
  \draw[white, opacity=0.5, line width=0.5pt] (0.2,0.2) rectangle (5.8,7.8);
  \node[white, font=\bfseries] at (3,5) {\rl{عنوان کتاب}};
  \node[white, font=\small] at (3,4.3) {\rl{زیرعنوان}};
  \node[matingold, font=\bfseries] at (3,1) {MatinBook};
\end{tikzpicture}
\end{latin}
\caption{نمای ساده‌شده‌ی جلد \lr{MatinBook}}
\label{fig:cover-intro}
\end{figure}

\mnote{در \lr{tikzpicture}، حتی برای متن کوتاه فارسی هم باید از \lr{\textbackslash rl\{...\}} استفاده کنید.}

شکل \cref{fig:cover-intro} نمای ساده‌شده‌ای از
جلد را نشان می‌دهد. فصل \cref{chap:cover} به‌طور
کامل به جلد می‌پردازد.

\section{مقایسه با سایر ابزارها}
\label{sec:introduction-comparison}

چندین کلاس و بسته‌ی \lr{LaTeX} برای نوشتن
کتاب‌های فارسی وجود دارد. در این بخش،
\lr{MatinBook} را با مهم‌ترین آن‌ها مقایسه
می‌کنیم.

\subsection{مقایسه با کلاس \lr{book} استاندارد}
\label{sec:introduction-vs-book}

کلاس \lr{book} استاندارد \lr{LaTeX}، پایه‌ی
\lr{MatinBook} است. ولی استفاده‌ی مستقیم از
آن برای کتاب فارسی، چالش‌های زیادی دارد:

\begin{itemize}
    \item \textbf{حروف‌چینی راست‌به‌چپ:} کلاس
    \lr{book} از RTL پشتیبانی نمی‌کند. باید
    \lr{xepersian} را دستی اضافه کنید.

    \item \textbf{فونت فارسی:} باید فونت فارسی
    را دستی تنظیم کنید.

    \item \textbf{محیط‌های علمی:} هیچ محیط علمی
    رنگ‌آمیزی‌شده‌ای ندارد.

    \item \textbf{کد و الگوریتم:} باید \lr{minted}
    و \lr{algorithm} را دستی پیکربندی کنید.

    \item \textbf{جلد:} هیچ طرح جلدی ندارد.
\end{itemize}

\lr{MatinBook} تمام این مشکلات را حل کرده است.

\subsection{مقایسه با \lr{XePersian} خالص}
\label{sec:introduction-vs-xepersian}

\lr{XePersian} یک بسته‌ی عالی برای حروف‌چینی
فارسی است، ولی یک بسته است، نه یک کلاس. اگر
از \lr{XePersian} خالص استفاده کنید، باید
همه‌ی موارد زیر را دستی پیکربندی کنید:

\begin{itemize}
    \item هندسه‌ی صفحه.
    \item سرصفحه و پاصفحه.
    \item محیط‌های علمی.
    \item سبک عناوین.
    \item کد و الگوریتم.
    \item جلد.
\end{itemize}

\mnote{تفاوت اصلی این است: \lr{XePersian} یک «موتور» است، \lr{MatinBook} یک «کتاب کامل» است.}

\lr{MatinBook} همه‌ی این‌ها را از پیش پیکربندی
شده ارائه می‌دهد.

\subsection{جدول مقایسه}
\label{sec:introduction-comparison-table}

جدول \cref{tab:comparison} مقایسه‌ی خلاصه‌ای
بین \lr{MatinBook}، کلاس \lr{book} استاندارد،
و \lr{XePersian} خالص را نشان می‌دهد.

\begin{matintable}{مقایسه‌ی \lr{MatinBook} با سایر ابزارها}{tab:comparison}
\begin{tblr}{
    width = \textwidth,
    colspec = {X[l] X[c] X[c] X[c]},
    hlines, vlines,
    colsep = 6pt,
    rowsep = 4pt,
    row{1} = {font=\bfseries},
}
قابلیت & \lr{MatinBook} & \lr{book} & \lr{XePersian} \\
حروف‌چینی RTL & \textcolor{matingreen}{بله} & \textcolor{matinred}{خیر} & \textcolor{matingreen}{بله} \\
فونت فارسی & \textcolor{matingreen}{خودکار} & \textcolor{matinred}{خیر} & \textcolor{matinorange}{دستی} \\
محیط‌های علمی & \textcolor{matingreen}{۱۰ نوع} & \textcolor{matinred}{خیر} & \textcolor{matinred}{خیر} \\
کد و الگوریتم & \textcolor{matingreen}{کامل} & \textcolor{matinred}{خیر} & \textcolor{matinred}{خیر} \\
نمودار & \textcolor{matingreen}{TikZ} & \textcolor{matinred}{خیر} & \textcolor{matinred}{خیر} \\
جدول مدرن & \textcolor{matingreen}{tblr} & \textcolor{matinorange}{tabular} & \textcolor{matinorange}{tabular} \\
نمایه‌ی فارسی & \textcolor{matingreen}{xindy} & \textcolor{matinred}{خیر} & \textcolor{matinorange}{دستی} \\
جلد & \textcolor{matingreen}{تمام‌رنگی} & \textcolor{matinred}{خیر} & \textcolor{matinred}{خیر} \\
ماژولار & \textcolor{matingreen}{۱۹ ماژول} & \textcolor{matinred}{خیر} & \textcolor{matinred}{خیر} \\
\end{tblr}
\end{matintable}

\section{معماری \lr{MatinBook}}
\label{sec:introduction-architecture}

\lr{MatinBook} یک کلاس ماژولار است. یعنی به‌جای
یک فایل بزرگ، از ۱۹ ماژول مستقل تشکیل شده است
که هر یک مسئولیت مشخصی دارند. این معماری،
نگهداری، توسعه و شخصی‌سازی کلاس را بسیار آسان‌تر
می‌کند.

\mnote{معماری ماژولار یعنی اگر فقط می‌خواهید یک ماژول را تغییر دهید، نیازی به دست زدن به بقیه‌ی کلاس ندارید.}

\subsection{ماژول‌های اصلی}
\label{sec:introduction-modules}

ماژول‌های \lr{MatinBook} به سه دسته تقسیم می‌شوند:

\textbf{۱. هسته (۲ ماژول):}
\begin{itemize}
    \item \file{mb-core.sty} — بارگذارنده‌ی
    بسته‌های پایه.
    \item \file{mb-utils.sty} — دستورات کمکی.
\end{itemize}

\textbf{۲. ماژول‌های کاربردی (۱۰ ماژول):}
\begin{itemize}
    \item \file{mb-math.sty} — ریاضیات.
    \item \file{mb-boxes.sty} — استایل‌های
    جعبه.
    \item \file{mb-theorem.sty} — محیط‌های
    قضیه.
    \item \file{mb-code.sty} — کد.
    \item \file{mb-layout.sty} — چیدمان.
    \item \file{mb-headings.sty} — عناوین.
    \item \file{mb-references.sty} — ارجاعات.
    \item \file{mb-graphics.sty} — گرافیک.
    \item \file{mb-index.sty} — نمایه.
    \item \file{mb-table.sty} — جداول.
\end{itemize}

\textbf{۳. تم و لوکال (۵ ماژول):}
\begin{itemize}
    \item \file{mb-theme-colors.sty} — پالت
    رنگی.
    \item \file{mb-theme-cover.sty} — جلد.
    \item \file{mb-theme-default.sty} —
    بارگذارنده‌ی تم.
    \item \file{fa-IR.sty} — لوکال فارسی.
    \item \file{en-US.sty} — لوکال انگلیسی.
\end{itemize}

\subsection{ترتیب بارگذاری}
\label{sec:introduction-loading-order}

ماژول‌ها در سه فاز بارگذاری می‌شوند:

\textbf{فاز ۱ (قبل از \lr{xepersian}):}
هسته، تم، و ماژول‌های کاربردی.

\textbf{فاز ۲:}
\lr{xepersian} — باید آخرین بسته باشد.

\textbf{فاز ۳ (بعد از \lr{xepersian}):}
فونت‌ها، تایپوگرافی، و لوکال.

\begin{remark}
    تغییر این ترتیب، باعث خطاهای عجیب و غریب
    می‌شود. اگر می‌خواهید ماژول جدیدی اضافه
    کنید، حتماً آن را در فاز مناسب قرار دهید.
    \label{rem:loading-order}
\end{remark}

معماری \lr{MatinBook} در \cref{chap:modules}
به‌طور کامل بررسی می‌شود.

\section{خطاهای رایج}
\label{sec:introduction-mistakes}

در این بخش، چند خطای رایج که کاربران جدید
\lr{MatinBook} ممکن است مرتکب شوند را
معرفی می‌کنیم:

\subsection{استفاده از \lr{pdfLaTeX} به‌جای \lr{XeLaTeX}}
\label{sec:introduction-mistake-engine}

\lr{MatinBook} فقط با موتور \lr{XeLaTeX} کار
می‌کند. اگر از \lr{pdfLaTeX} استفاده کنید،
خطای زیر را دریافت می‌کنید:

\begin{latin}
\begin{minted}{text}
! Package matinbook Error:
  MatinBook requires XeLaTeX.
\end{minted}
\end{latin}

\mnote{\lr{XeLaTeX} برای پشتیبانی از فونت‌های \lr{OpenType} و حروف‌چینی راست‌به‌چپ ضروری است.}

\textbf{راه‌حل:} همیشه از \lr{XeLaTeX} استفاده
کنید.

\subsection{فراموش کردن \lr{-shell-escape}}
\label{sec:introduction-mistake-shellescape}

\lr{MatinBook} از \lr{minted} و \lr{xindy}
استفاده می‌کند که نیاز به \lr{-shell-escape}
دارند. اگر آن را فراموش کنید، خطاهای زیر
را دریافت می‌کنید:

\begin{itemize}
    \item \lr{minted}: عدم توانایی در اجرای
    \lr{Pygments}.
    \item \lr{xindy}: عدم توانایی در ساخت نمایه.
\end{itemize}

\textbf{راه‌حل:} همیشه از
\lr{xelatex -shell-escape} استفاده کنید.

\subsection{فراموش کردن \lr{biber}}
\label{sec:introduction-mistake-biber}

اگر از \lr{biblatex} و \lr{biber} استفاده
می‌کنید، باید \lr{biber} را دستی اجرا کنید:

\begin{latin}
\begin{minted}{bash}
xelatex -shell-escape main.tex
biber main
xelatex -shell-escape main.tex
xelatex -shell-escape main.tex
\end{minted}
\end{latin}

\mnote{این چهار مرحله را می‌توانید با \lr{latexmk} خودکار کنید: \lr{latexmk -xelatex -shell-escape main.tex}.}

\textbf{راه‌حل:} این چرخه‌ی چهار مرحله‌ای را
همیشه برای کتاب‌های حاوی منابع اجرا کنید.

\subsection{وارد کردن متن فارسی در بلوک کد}
\label{sec:introduction-mistake-persian-code}

بلوک‌های \env{minted} از فونت لاتین استفاده
می‌کنند، پس **نمی‌توانند متن فارسی را مستقیماً
نمایش دهند**. اگر متن فارسی در بلوک کد بنویسید،
نتیجه به‌صورت نامفهوم نمایش داده می‌شود.

\textbf{راه‌حل:} بسته به زبان، از یکی از دو
دستور زیر استفاده کنید:

\begin{itemize}
    \item \textbf{پایتون:} از \cmd{pcc} استفاده
    کنید.
    \item \textbf{سایر زبان‌ها:} از
    \cmd{pc} داخل \lr{|...|} استفاده کنید.
\end{itemize}

\begin{latin}
\begin{minted}{python}
x = 5  |\pcc{مقدار متغیر}|
y = 10  |\pcc{مقدار دیگر}|
\end{minted}
\end{latin}

\mnote{دستور \lr{\textbackslash pcc} توسط \lr{mb-code.sty} تعریف شده و علامت \lr{\#} را به‌طور خودکار اضافه می‌کند.}

خطاهای رایج در \cref{app:troubleshooting}
به‌طور کامل بررسی می‌شوند.

\section{تمرین‌ها}
\label{sec:introduction-exercises}

\begin{exercise}
    سه قابلیت \lr{MatinBook} که برای کار شما
    مهم‌تر هستند را فهرست کنید و توضیح دهید
    چرا.
    \label{exc:introduction-features}
\end{exercise}

\begin{exercise}
    \lr{MatinBook} را با یکی از ابزارهای
    حروف‌چینی که قبلاً استفاده کرده‌اید
    مقایسه کنید. سه مزیت و یک عیب
    \lr{MatinBook} را بنویسید.
    \label{exc:introduction-compare}
\end{exercise}

\begin{exercise}
    در فصل‌های بعدی، چه موضوعاتی را می‌خواهید
    اول یاد بگیرید؟ یک فهرست اولویت‌بندی‌شده
    تهیه کنید.
    \label{exc:introduction-priority}
\end{exercise}

\section{خلاصه}
\label{sec:introduction-summary}

در این فصل، با \lr{MatinBook} آشنا شدیم:

\begin{itemize}
    \item \lr{MatinBook} یک کلاس \lr{LaTeX}
    است که برای کتاب‌های فنی فارسی طراحی
    شده است.

    \item این کلاس از موتور \lr{XeLaTeX}
    استفاده می‌کند و از حروف‌چینی راست‌به‌چپ
    کامل پشتیبانی می‌کند.

    \item قابلیت‌های کلیدی شامل محیط‌های علمی
    رنگ‌آمیزی‌شده، ریاضیات، کد، نمودار، جدول،
    جلد حرفه‌ای و نمایه‌ی فارسی است.

    \item \lr{MatinBook} از ۱۹ ماژول مستقل
    تشکیل شده که در سه فاز بارگذاری
    می‌شوند.

    \item خطاهای رایج شامل استفاده از
    \lr{pdfLaTeX}، فراموش کردن
    \lr{-shell-escape}، فراموش کردن
    \lr{biber} و وارد کردن متن فارسی در
    بلوک کد است.
\end{itemize}

در فصل بعدی (\cref{chap:installation})،
به‌طور گام‌به‌گام یاد می‌گیرید که
\lr{MatinBook} را نصب کنید و اولین کتاب
خود را بسازید.